\documentclass{article}%
\usepackage[T1]{fontenc}%
\usepackage[utf8]{inputenc}%
\usepackage{lmodern}%
\usepackage{textcomp}%
\usepackage{lastpage}%
\usepackage{geometry}%
\geometry{margin=2.2cm}%
%
\usepackage{amsmath}%
\usepackage{amssymb}%
\usepackage{booktabs}%
\usepackage{longtable}%
\usepackage{array}%
\usepackage{hyperref}%
\hypersetup{hidelinks}%
\title{ITA 2027 {-} Matemática (claude\_sonnet\_5\_5) x Poliedro\\[4pt]\large Comparação com o gabarito Poliedro Resolve (ITA 2027, 2ª fase, dia 1 {-} Matemática)\thanks{Gerado com gpt-resolve: \url{https://github.com/lgabs/gpt-resolve}}}%
\date{09/11/2026}%
%
\begin{document}%
\normalsize%
\maketitle%
\section*{Resumo}\textbf{Nota média: 10,0/10} (12 de 12 itens com a resolução correta; cada item vale o mesmo dentro da questão).\par\medskip%
\begin{longtable}{@{}c c p{0.78\textwidth}@{}}\toprule Questão & Nota & Resumo \\ \midrule
1 & 10,0 & A resolução e o gabarito concordam: as PAs têm razão ímpar com |r| $\leq$ 9 e a$_{1}$ = (87 - 9r)/2, dando 5 crescentes, (39,\ldots{},48), (30,\ldots{},57), (21,\ldots{},66), (12,\ldots{},75) e (3,\ldots{},84), mais as 5 decrescentes correspondentes, 10 no total. Confirmei por força bruta em python que existem exatamente essas 10. A resolução é mais completa que o gabarito, porque justifica a paridade, o limite r $\leq$ 9 e a simetria por inversão. O único reparo, menor, é que ela deixa a contagem final condicional ('5, ou 10 com as decrescentes') em vez de afirmar 10. Ainda assim, lista as 10 explicitamente. Veredito: correta, de acordo com o gabarito. \\
2 & 10,0 & A resolução está correta e bate com o gabarito do Poliedro: S = (log2(($\surd$5-1)/2), 1), que também pode ser escrito (-log2((1+$\surd$5)/2), 1) $\approx$ (-0,694; 1). Refiz a solução analiticamente e conferi com uma varredura numérica: o intervalo e a exclusão dos dois extremos estão confirmados. Só há detalhes menores: um sinal '<' trocado na resolução do gabarito (Caso 2), que parece erro de digitação, e uma justificativa mal formulada no Passo 1 da resolução. Nenhum dos dois afeta a resposta. \\
3 & 10,0 & A resolução está correta e bate com o gabarito do Poliedro nos dois itens: r = 1/16 e lado ? = $\surd$(4-2$\surd$2)/2 = $\surd$((2-$\surd$2)/2) $\approx$ 0,541. O caminho é o mesmo do gabarito: De Moivre, octógono regular de raio r\^{}\{1/8\}, área por 8 triângulos e lado pela lei dos cossenos. A resolução ainda confere o resultado com a fórmula da área do octógono pelo lado. A transcrição do enunciado confere com o PDF original. Conferi raízes, área e lado em Python. \\
4 & 10,0 & A resolução está correta e coincide com o gabarito do Poliedro: a = {-}2, e as raízes são {-}1 (dupla), 2, $\pm$i e 3$\pm$2$\surd$2. Refiz a divisão, a fatoração completa e o cálculo de p, p' e s nas raízes de q com aritmética exata. A resolução é até mais completa que o gabarito, porque verifica de forma independente as três condições p(2) = 0, p({-}1) = 0 e p'({-}1) = 0. Não encontrei erro. \\
5 & 10,0 & A resolução está certa e bate com o gabarito do Poliedro: x\_C = 5$\surd$3 - 9. Os dois seguem o mesmo caminho: C fica na interseção da elipse dada com a elipse de focos A e B e eixo maior 6, o que leva a x$^2$+18x+6=0, e a raiz -9-5$\surd$3 é descartada porque cai fora de |x|$\leq$$\surd$3. Conferi todas as contas em Python, e o perímetro dá 10. Não há erro de matemática em nenhum dos dois, só falhas de formatação na transcrição do gabarito. \\
6 & 10,0 & A resolução está correta e coincide com o gabarito do Poliedro: são 235 sêxtuplas. A redução do sistema dá x2 = 6, x1 = 8 - 2x6 e x3+x4+x5 = 16 - x6, com x6 $\in$ \{1,2,3\}. A contagem por estrelas e barras é 91 + 78 + 66 = 235. Refiz as contas e testei por força bruta em Python, que também deu 235. O enunciado transcrito confere com o PDF original. \\
7 & 10,0 & A resolução e o gabarito estão certos e chegam à mesma resposta, n $\in$ \{-6, -4, 6, 8\}. O método é o mesmo nos dois: dividir os polinômios (resto 63), exigir que n$^2$-2n-27 divida 63 e testar quais divisores tornam 28+d um quadrado perfeito (d = -3 e d = 21). A resolução está correta e completa e ainda confere os quocientes. O gabarito tem um erro de digitação sem importância na expressão do discriminante. \\
8 & 10,0 & A resolução está certa nos dois itens e coincide com o gabarito do Poliedro. No item a, ambos dão 0 < x $\leq$ 1 ou 9 $\leq$ x < 10. No item b, ambos dão lados $\surd$10 e 3$\surd$10. Refiz as contas e chego ao mesmo resultado, que também conferi numericamente. A resolução é até mais completa: mostra que y fica em (0,6) e verifica o ângulo reto e a área. O gabarito tem só um erro de digitação ('x-5y-3=0' no lugar de (x-5)(y-3)=0), que não afeta a resposta. \\
9 & 10,0 & A resolução está certa e bate com o gabarito do Poliedro: a área da seção é 5(2+$\surd$3) = 10+5$\surd$3 $\approx$ 18,66. Os dois fazem o mesmo caminho. A seção é um trapézio de bases 8 (diagonal maior AD) e 2 (base média PQ = EF/2). A altura é tg75$^\circ$ = 2+$\surd$3, ou seja, metade do apótema lateral VM = 2(2+$\surd$3). Conferi com uma conta em coordenadas no Python: o ângulo no vértice dá 30$^\circ$, a seção tem os vértices ($\pm$4,0,0) e ($\pm$1,-$\surd$3,h/2), e a área é 18,660254\ldots{}, igual a 5(2+$\surd$3). A única falha é no texto do gabarito, onde a fórmula final saiu truncada. O valor final do gabarito está certo. \\
10 & 10,0 & A resolução está correta e concorda com o gabarito do Poliedro: k = 3. No 1º quadrante, f percorre ($\surd$2-1, 2-$\surd$2{]}, intervalo sem inteiros. No 3º quadrante, percorre ($\surd$2+1, 2+$\surd$2{]}, que contém só o inteiro 3. Confirmei isso numericamente e conferi que o x explícito dado na resolução ($\approx$ 3,429 rad, no 3º quadrante) dá f(x) = 3. A resolução é até mais cuidadosa que o gabarito: usa os módulos corretamente, justifica por continuidade que todos os valores do intervalo são atingidos e trata o caso dos pontos sobre os eixos. O gabarito só tem pequenas imprecisões de redação, como a falta dos módulos numa linha intermediária e a mistura de desigualdades estritas e não estritas, que não afetam a resposta. \\
\bottomrule\end{longtable}%
\section*{Questão 1}%
\href{https://poliedroresolve.sistemapoliedro.com.br/vestibulares/ita/2027/ita-2-fase-dia-1-20-10-2026-2-fase-dia-1/dissertativa-1-matematica-2-fase-dia-1-ita-2027}{Resolução do gabarito}\par\smallskip%
\begin{longtable}{@{}p{0.06\textwidth} p{0.33\textwidth} p{0.33\textwidth} p{0.1\textwidth} p{0.08\textwidth}@{}}\toprule Item & Resolução & Gabarito & Veredito & Certo \\ \midrule
Q1 (única) & r ímpar, 1 $\leq$ r $\leq$ 9, a$_{1}$ = (87 - 9r)/2: (39,\ldots{},48) r=1; (30,\ldots{},57) r=3; (21,\ldots{},66) r=5; (12,\ldots{},75) r=7; (3,\ldots{},84) r=9; mais as 5 invertidas (r = -1, -3, -5, -7, -9). Ao todo 5 PAs crescentes, ou 10 contando as decrescentes. & 10 PAs: as crescentes (39,\ldots{},48) r=1; (30,\ldots{},57) r=3; (21,\ldots{},66) r=5; (12,\ldots{},75) r=7; (3,\ldots{},84) r=9, e as 5 decrescentes correspondentes (r = -1, -3, -5, -7, -9). & igual & ambos \\
\bottomrule\end{longtable}%
\begin{itemize}%
\item%
\textbf{Q1 (única))} Reli o enunciado no PDF original (pág. 1): 'Descreva todas as progressões aritméticas de 10 números naturais cuja soma resulte em 435', com $\mathbb{N}$ = \{1,2,3,\ldots{}\}. A transcrição está correta. A condição S$_{1}$$_{0}$ = 5(2a$_{1}$+9r) = 435 dá 2a$_{1}$ + 9r = 87. Daí r é ímpar, e a exigência de que o menor termo seja $\geq$ 1 dá |r| $\leq$ 9. Fiz uma busca por força bruta em python (a$_{1}$ de 1 a 435, r de -100 a 100) e encontrei exatamente 10 soluções: (a$_{1}$,r) = (3,9), (12,7), (21,5), (30,3), (39,1), (48,-1), (57,-3), (66,-5), (75,-7), (84,-9). São exatamente as 10 da resolução e as 10 do gabarito. A resolução descreve as 10 sequências explicitamente (as 5 crescentes e as 5 invertidas, com razões e primeiros termos). Ela só deixa a contagem em aberto ('5 crescentes, ou 10 contando as decrescentes'), enquanto o gabarito afirma 10. Como as sequências invertidas são PAs distintas e válidas, o certo é contar 10, e a resolução de fato as inclui. Por isso o conteúdo coincide.%
\end{itemize}%
\paragraph{Pontos de dúvida}%
\begin{itemize}%
\item%
\textit{(solution)} A resolução fecha com uma contagem condicional: '5 PAs crescentes, ou 10 contando as decrescentes'. Como (48,47,\ldots{},39) é uma PA diferente de (39,\ldots{},48), o enunciado ('todas as PAs') pede 10, que é o que o gabarito afirma. As 10 estão descritas, mas um corretor rigoroso pode preferir uma afirmação firme ('são 10 PAs') a deixar a contagem a cargo do leitor.%
\item%
\textit{(key)} A resolução do Poliedro tem problemas de notação. Faltam parênteses em 'S\_10 = (a\_1 + a\_10 $\cdot$ 10)/2' e em '(2a\_1 + 9r $\cdot$ 10)/2'. Além disso, a lista das decrescentes aparece concatenada, sem separadores ('\ldots{}40,3957,54,\ldots{}'), o que parece ser artefato de extração. Nada disso afeta o resultado.%
\item%
\textit{(key)} O gabarito não justifica por que a lista de razões para em r = 9. Falta a condição a$_{1}$ = (87-9r)/2 $\geq$ 1 $\Rightarrow$ r $\leq$ 85/9; com r = 11 daria a$_{1}$ = -6. Também não justifica a redução das razões negativas à inversão das crescentes. A resolução conferida faz as duas justificativas de forma explícita.%
\end{itemize}%
\section*{Questão 2}%
\href{https://poliedroresolve.sistemapoliedro.com.br/vestibulares/ita/2027/ita-2-fase-dia-1-20-10-2026-2-fase-dia-1/dissertativa-2-matematica-2-fase-dia-1-ita-2027}{Resolução do gabarito}\par\smallskip%
\begin{longtable}{@{}p{0.06\textwidth} p{0.33\textwidth} p{0.33\textwidth} p{0.1\textwidth} p{0.08\textwidth}@{}}\toprule Item & Resolução & Gabarito & Veredito & Certo \\ \midrule
Questão 2 (conjunto solução de |log2(2\^{}x+2)| > |2x|) & x $\in$ (log2(($\surd$5-1)/2), 1) = (-log2((1+$\surd$5)/2), 1) $\approx$ (-0,694; 1) & S = \{ x $\in$ $\mathbb{R}$ | log2(($\surd$5 - 1)/2) < x < 1 \} & igual & ambos \\
\bottomrule\end{longtable}%
\begin{itemize}%
\item%
\textbf{Questão 2 (conjunto solução de |log2(2\^{}x+2)| > |2x|))} Os dois dão o mesmo intervalo aberto, e a resolução também escreve o extremo inferior na forma equivalente -log2($\varphi$), com $\varphi$ = (1+$\surd$5)/2. Refiz a questão. Como 2\^{}x + 2 > 2, o log é maior que 1 e o módulo da esquerda pode ser retirado. Caso x $\geq$ 0: t$^2$ - t - 2 < 0 dá 2\^{}x < 2, ou seja, 0 $\leq$ x < 1. Caso x < 0: t$^3$ + 2t$^2$ - 1 = (t+1)(t$^2$ + t - 1) > 0 dá t > ($\surd$5-1)/2, ou seja, log2(($\surd$5-1)/2) < x < 0. A união é (log2(($\surd$5-1)/2), 1). Uma varredura numérica em {[}-5, 5{]} com passo 10$^{-}$$^{4}$ deu exatamente {[}-0,6942; 0,9999{]}, sem buracos, e log2(($\surd$5-1)/2) $\approx$ -0,69424. Nos dois extremos vale a igualdade, então ambos ficam de fora, como os dois textos afirmam. Os valores de teste da resolução (x = -0,69, -0,70, 0,5 e 2) também conferem.%
\end{itemize}%
\paragraph{Pontos de dúvida}%
\begin{itemize}%
\item%
\textit{(key)} Na resolução do gabarito, o Caso 2 aparece como 'log\_2(2\^{}x+2) < -2x $\Rightarrow$ 2\^{}x + 2 > 2\^{}(-2x)'. O sinal '<' está invertido: deveria ser '>'. Parece erro de digitação ou de extração do texto, porque a linha seguinte e a conclusão usam o sentido certo, e o resultado final não muda.%
\item%
\textit{(solution)} No Passo 1, a resolução justifica a passagem log2(2\^{}x+2) > |2x| $\Rightarrow$ 2\^{}x + 2 > 2\^{}|2x| dizendo 'Como |2x| $\geq$ 0 e log2 é crescente'. O motivo certo é que 2\^{}($\cdot$) é crescente; o fato de |2x| $\geq$ 0 não tem papel aí. É só imprecisão de redação, sem efeito na conta.%
\end{itemize}%
\section*{Questão 3}%
\href{https://poliedroresolve.sistemapoliedro.com.br/vestibulares/ita/2027/ita-2-fase-dia-1-20-10-2026-2-fase-dia-1/dissertativa-3-matematica-2-fase-dia-1-ita-2027}{Resolução do gabarito}\par\smallskip%
\begin{longtable}{@{}p{0.06\textwidth} p{0.33\textwidth} p{0.33\textwidth} p{0.1\textwidth} p{0.08\textwidth}@{}}\toprule Item & Resolução & Gabarito & Veredito & Certo \\ \midrule
lado do polígono & ? = $\surd$((2-$\surd$2)/2) = $\surd$(4-2$\surd$2)/2 $\approx$ 0,541 & l = $\surd$(4 - 2$\surd$2)/2 = $\surd$((2 - $\surd$2)/2) & igual & ambos \\
valor de r & r = 1/16 & r = 1/16 & igual & ambos \\
\bottomrule\end{longtable}%
\begin{itemize}%
\item%
\textbf{lado do polígono)} Mesma expressão, escrita nas mesmas duas formas. Conferi numericamente: com r = 1/16, a distância entre raízes vizinhas dá 0,541196, igual a $\surd$(4-2$\surd$2)/2.%
\item%
\textbf{valor de r)} z\^{}8 = ir dá um octógono regular de raio R = r\^{}\{1/8\}. A área é 2$\surd$2$\cdot$R$^2$ = $\surd$2, logo R$^2$ = 1/2 e r = R\^{}8 = 1/16. Conferi numericamente que as oito raízes satisfazem r + iz\^{}8 = 0 e que o polígono tem área $\surd$2.%
\end{itemize}%
\section*{Questão 4}%
\href{https://poliedroresolve.sistemapoliedro.com.br/vestibulares/ita/2027/ita-2-fase-dia-1-20-10-2026-2-fase-dia-1/dissertativa-4-matematica-2-fase-dia-1-ita-2027}{Resolução do gabarito}\par\smallskip%
\begin{longtable}{@{}p{0.06\textwidth} p{0.33\textwidth} p{0.33\textwidth} p{0.1\textwidth} p{0.08\textwidth}@{}}\toprule Item & Resolução & Gabarito & Veredito & Certo \\ \midrule
Q4 (valor de a e raízes de p) & a = {-}2; raízes: {-}1 (dupla), 2, i, {-}i, 3+2$\surd$2, 3{-}2$\surd$2 & a = {-}2; raízes: {-}1 (dupla), 2, i, {-}i, 3+2$\surd$2, 3{-}2$\surd$2 & igual & ambos \\
\bottomrule\end{longtable}%
\begin{itemize}%
\item%
\textbf{Q4 (valor de a e raízes de p))} Conferi com aritmética exata. O resto de p por q, em função de a, sai proporcional a (a+2): (10+5a)x$^2$ + (52+26a)x + (26+13a). Ele só se anula para a = {-}2. Com a = {-}2, o quociente é x$^{4}${-}6x$^3$+2x$^2${-}6x+1, com resto zero. O produto (x+1)$^2$(x{-}2)(x$^2$+1)(x$^2${-}6x+1) reproduz exatamente p(x) = x$^{7}${-}6x$^{6}${-}x$^{5}$+10x$^{4}$+7x$^3$+14x$^2$+9x{-}2. Conferi também a resolução: p(2) = 170+85a, p({-}1) = {-}16{-}8a e p'({-}1) = 32+16a, todas corretas. Como s({-}1) = 16 e s(2) = {-}35, as multiplicidades são 2 para {-}1 e 1 para 2. O enunciado do PDF (pág. 5 do caderno) bate com a transcrição.%
\end{itemize}%
\paragraph{Pontos de dúvida}%
\begin{itemize}%
\item%
\textit{(key)} O gabarito tira a = {-}2 só de p({-}1) = 0, sem usar p(2) = 0 nem a multiplicidade dupla de {-}1. A compatibilidade das condições só fica garantida depois, porque a divisão dá resto nulo. O argumento está completo, mas seria mais claro dizer que a divisão exata confirma a hipótese.%
\end{itemize}%
\section*{Questão 5}%
\href{https://poliedroresolve.sistemapoliedro.com.br/vestibulares/ita/2027/ita-2-fase-dia-1-20-10-2026-2-fase-dia-1/dissertativa-5-matematica-2-fase-dia-1-ita-2027}{Resolução do gabarito}\par\smallskip%
\begin{longtable}{@{}p{0.06\textwidth} p{0.33\textwidth} p{0.33\textwidth} p{0.1\textwidth} p{0.08\textwidth}@{}}\toprule Item & Resolução & Gabarito & Veredito & Certo \\ \midrule
abscissa de C & x\_C = 5$\surd$3 - 9 ($\approx$ -0,34) & x = 5$\surd$3 - 9 & igual & ambos \\
\bottomrule\end{longtable}%
\begin{itemize}%
\item%
\textbf{abscissa de C)} Mesmo valor. Refiz a conta: com x=1 na elipse 2x$^2$+y$^2$=6, y=$\pm$2, logo AB=4 e CA+CB=6. C está então na elipse de focos A e B, (x-1)$^2$/5 + y$^2$/9 = 1, que equivale a 9x$^2$-18x+5y$^2$=36. Fazendo 5$\cdot$(2x$^2$+y$^2$=6) - (9x$^2$-18x+5y$^2$=36), sai x$^2$+18x+6=0, com raízes x=-9$\pm$5$\surd$3. Como |x|$\leq$$\surd$3, só fica 5$\surd$3-9$\approx$-0,3397. Conferi em Python: y$\approx$2,4019, CA$\approx$1,3987, CB$\approx$4,6013, perímetro = 10,000. O valor x$^2$=156-90$\surd$3 citado pela resolução também está certo. Como x$\neq$1, o triângulo não é degenerado.%
\end{itemize}%
\paragraph{Pontos de dúvida}%
\begin{itemize}%
\item%
\textit{(key)} A transcrição da resolução do gabarito tem erros de digitação que não afetam o resultado: "A = 1,2" e "centro 1,0" estão sem parênteses, e a fórmula aparece como "(x{-}1\^{}2)/5" no lugar de (x-1)$^2$/5. A forma expandida que o gabarito usa, 9x$^2$-18x+5y$^2$=36, está correta.%
\item%
\textit{(solution)} A resolução fixa A=(1,2) e B=(1,-2) sem perda de generalidade. Trocar os rótulos não muda a soma CA+CB nem a abscissa de C, então está tudo certo. Ela também observa que há dois pontos C possíveis, (x, $\pm$y), com a mesma abscissa, e isso é coerente com o enunciado, que pede apenas a abscissa.%
\end{itemize}%
\section*{Questão 6}%
\href{https://poliedroresolve.sistemapoliedro.com.br/vestibulares/ita/2027/ita-2-fase-dia-1-20-10-2026-2-fase-dia-1/dissertativa-6-matematica-2-fase-dia-1-ita-2027}{Resolução do gabarito}\par\smallskip%
\begin{longtable}{@{}p{0.06\textwidth} p{0.33\textwidth} p{0.33\textwidth} p{0.1\textwidth} p{0.08\textwidth}@{}}\toprule Item & Resolução & Gabarito & Veredito & Certo \\ \midrule
Questão 6 (número de sêxtuplas) & 235 sêxtuplas ordenadas & 235 sêxtuplas & igual & ambos \\
\bottomrule\end{longtable}%
\begin{itemize}%
\item%
\textbf{Questão 6 (número de sêxtuplas))} Conferi o enunciado no PDF original (pág. 6) e ele bate com a transcrição. A redução feita é válida: com S = x3+x4+x5, sai x2 = 6, x1 = 8 - 2x6 e S = 16 - x6. Como x1 $\geq$ 1, temos x6 $\leq$ 3, logo x6 $\in$ \{1,2,3\}. Pelo método das estrelas e barras, o número de ternas positivas com soma S é C(S-1,2), o que dá C(14,2)+C(13,2)+C(12,2) = 91+78+66 = 235. O gabarito chega ao mesmo resultado com outra notação: P\_14\^{}\{12,2\} = 14!/(12!$\cdot$2!) = 91, e assim por diante. Também testei todas as sêxtuplas por força bruta em Python, com cada xi $\geq$ 1, e o total foi 235.%
\end{itemize}%
\paragraph{Pontos de dúvida}%
\begin{itemize}%
\item%
\textit{(key)} O gabarito afirma que 1 $\leq$ x6 $\leq$ 3 sem justificar. O limite vem de x1 = 8 - 2x6 $\geq$ 1, junto com o fato de x6 ser inteiro. A resolução conferida faz essa justificativa. Isso afeta só a redação, não o resultado.%
\end{itemize}%
\section*{Questão 7}%
\href{https://poliedroresolve.sistemapoliedro.com.br/vestibulares/ita/2027/ita-2-fase-dia-1-20-10-2026-2-fase-dia-1/dissertativa-7-matematica-2-fase-dia-1-ita-2027}{Resolução do gabarito}\par\smallskip%
\begin{longtable}{@{}p{0.06\textwidth} p{0.33\textwidth} p{0.33\textwidth} p{0.1\textwidth} p{0.08\textwidth}@{}}\toprule Item & Resolução & Gabarito & Veredito & Certo \\ \midrule
Questão 7 (valores inteiros de n) & n $\in$ \{-6, -4, 6, 8\} (quocientes -10, -30, -10 e 18) & n $\in$ \{-6, -4, 6, 8\} & igual & ambos \\
\bottomrule\end{longtable}%
\begin{itemize}%
\item%
\textbf{Questão 7 (valores inteiros de n))} A transcrição confere com o enunciado original (p. 8 do PDF). A divisão está certa: (2n-1)(n$^2$-2n-27) = 2n$^3$-5n$^2$-52n+27, então o resto é 90-27 = 63 e o quociente é E = 2n-1 + 63/(n$^2$-2n-27). O denominador nunca se anula para n inteiro, porque (n-1)$^2$ = 28 não tem solução inteira. Como 28+d precisa ser quadrado perfeito, só sobram d = -3, que dá (n-1)$^2$ = 25 e n $\in$ \{-4, 6\}, e d = 21, que dá (n-1)$^2$ = 49 e n $\in$ \{-6, 8\}. Confirmei por força bruta em python, para |n| $\leq$ 100000, que estes são os únicos valores. Os quocientes saem -10, -30, -10 e 18, iguais aos da resolução.%
\end{itemize}%
\paragraph{Pontos de dúvida}%
\begin{itemize}%
\item%
\textit{(key)} Há erro de digitação no discriminante do gabarito. Para n$^2$ - 2n - (k+27) = 0, o certo é $\Delta$ = 4 + 4(k+27) = 4(k+28), mas o gabarito escreve '$\Delta$ = 4 + 4k + 27' e 'm$^2$ = 4k + 28'. A conclusão (k+28 tem de ser quadrado perfeito) e a resposta final não mudam.%
\item%
\textit{(solution)} A resolução cita 'uma busca computacional em -2000 $\leq$ n $\leq$ 2000', o que não vale como argumento numa prova escrita. A demonstração algébrica já está completa sem isso: D(n) divide 63 e 28+D(n) tem de ser quadrado perfeito. Na versão entregue, convém tirar essa frase ou deixar claro que é só uma conferência.%
\end{itemize}%
\section*{Questão 8}%
\href{https://poliedroresolve.sistemapoliedro.com.br/vestibulares/ita/2027/ita-2-fase-dia-1-20-10-2026-2-fase-dia-1/dissertativa-8-matematica-2-fase-dia-1-ita-2027}{Resolução do gabarito}\par\smallskip%
\begin{longtable}{@{}p{0.06\textwidth} p{0.33\textwidth} p{0.33\textwidth} p{0.1\textwidth} p{0.08\textwidth}@{}}\toprule Item & Resolução & Gabarito & Veredito & Certo \\ \midrule
a & 0 < x $\leq$ 1 ou 9 $\leq$ x < 10 & 0 < x $\leq$ 1 ou 9 $\leq$ x < 10 & igual & ambos \\
b & Lados $\surd$10 e 3$\surd$10 (para x = 1 ou x = 9, com y = 3) & Lados 3$\surd$10 e $\surd$10 & igual & ambos \\
\bottomrule\end{longtable}%
\begin{itemize}%
\item%
\textbf{a)} Os dois chegam à mesma condição, y(6-y) = x(10-x), com y = BF. A resolução usa coordenadas e produto escalar; o gabarito usa semelhança entre EBF e HAE. Para existir y em (0,6), é preciso x(10-x) $\leq$ 9, isto é, (x-1)(x-9) $\geq$ 0, que é o mesmo $\Delta$ $\geq$ 0 do gabarito. Restringindo a 0<x<10, sai o intervalo da resposta. A resolução ainda confere que y = 3 $\pm$ $\surd$((x-1)(x-9)) fica em (0,6), o que o gabarito não faz. Varri x de 0,001 a 9,999 e o resultado bateu com (0,1{]} ? {[}9,10).%
\item%
\textbf{b)} A área do retângulo inscrito é 60 - 6x - 10y + 2xy. Igualando a 30, obtém{-}se (x-5)(y-3) = 0. x = 5 é descartado pelo item a, então y = 3 e, daí, x(10-x) = 9, ou seja, x = 1 ou x = 9. Para x = 1, EF = $\surd$(81+9) = 3$\surd$10 e EH = $\surd$(1+9) = $\surd$10, com produto 30. Para x = 9 os valores apenas se invertem. Conferi numericamente: a área dá 30 e a relação (1) dá resíduo 0.%
\end{itemize}%
\paragraph{Pontos de dúvida}%
\begin{itemize}%
\item%
\textit{(key)} No item b, o gabarito escreve 'x-5y-3 = 0'. É erro de digitação de (x-5)(y-3) = 0: a sequência ('x = 5 é impossível, logo y = 3') confirma a fatoração correta. A resposta final não muda.%
\item%
\textit{(statement)} Os dois leem 'inscrito' como um vértice de EFGH em cada lado de ABCD, com G em CD e H em DA. É a leitura padrão e o enunciado não a declara. A resolução deixa essa hipótese explícita e mostra que G = (10-x, 6) cai no interior de CD.%
\end{itemize}%
\section*{Questão 9}%
\href{https://poliedroresolve.sistemapoliedro.com.br/vestibulares/ita/2027/ita-2-fase-dia-1-20-10-2026-2-fase-dia-1/dissertativa-9-matematica-2-fase-dia-1-ita-2027}{Resolução do gabarito}\par\smallskip%
\begin{longtable}{@{}p{0.06\textwidth} p{0.33\textwidth} p{0.33\textwidth} p{0.1\textwidth} p{0.08\textwidth}@{}}\toprule Item & Resolução & Gabarito & Veredito & Certo \\ \midrule
Área da seção & 5(2+$\surd$3) = 10+5$\surd$3 $\approx$ 18,66 (trapézio isósceles de bases 8 e 2, altura 2+$\surd$3) & 5(2+$\surd$3) & igual & ambos \\
\bottomrule\end{longtable}%
\begin{itemize}%
\item%
\textbf{Área da seção)} Refiz a conta. A aresta lateral é ?$^2$ = 16/(2-$\surd$3) = 16(2+$\surd$3), a altura da pirâmide é h$^2$ = 16+16$\surd$3 e o apótema lateral é VM = 2(2+$\surd$3) = 2$\cdot$tg75$^\circ$. O plano paralelo à face VBC que passa por AD corta as arestas VE e VF nos pontos médios. Isso dá um trapézio de bases 8 e 2 com altura OT = VM/2 = 2+$\surd$3. Área = (8+2)/2$\cdot$(2+$\surd$3) = 5(2+$\surd$3). Uma verificação numérica com vetores (interseção do plano com as arestas e área do polígono) deu 18,660254, que coincide com 5(2+$\surd$3).%
\end{itemize}%
\paragraph{Pontos de dúvida}%
\begin{itemize}%
\item%
\textit{(key)} No texto da resolução do gabarito, a fórmula final saiu truncada: "(AD + PQOT)/2 = (8 + 2 $\cdot$ tg75$^\circ$)/2 = 52 + $\surd$(3)". O certo seria (AD+PQ)$\cdot$OT/2 = (8+2)$\cdot$tg75$^\circ$/2 = 5(2+$\surd$3). Parece só perda de formatação na transcrição, porque o valor final 5(2+$\surd$3) está correto.%
\item%
\textit{(statement)} O enunciado não diz a qual face lateral a seção é paralela. Só as duas faces cuja aresta da base é paralela à diagonal escolhida tornam a seção possível, e por simetria elas dão seções congruentes. A resolução justifica isso explicitamente; o gabarito assume sem comentar. Não muda a resposta.%
\end{itemize}%
\section*{Questão 10}%
\href{https://poliedroresolve.sistemapoliedro.com.br/vestibulares/ita/2027/ita-2-fase-dia-1-20-10-2026-2-fase-dia-1/dissertativa-10-matematica-2-fase-dia-1-ita-2027}{Resolução do gabarito}\par\smallskip%
\begin{longtable}{@{}p{0.06\textwidth} p{0.33\textwidth} p{0.33\textwidth} p{0.1\textwidth} p{0.08\textwidth}@{}}\toprule Item & Resolução & Gabarito & Veredito & Certo \\ \midrule
único & k = 3 (no 1º quadrante f $\in$ ($\surd$2-1, 2-$\surd$2{]}; no 3º quadrante f $\in$ ($\surd$2+1, 2+$\surd$2{]}) & k = 3 & igual & ambos \\
\bottomrule\end{longtable}%
\begin{itemize}%
\item%
\textbf{único)} Conferi o enunciado no PDF original (pág. 10) e ele coincide com a transcrição. Simplificando, f(x) = $\surd$2(|cos x|+|sen x|) - (sen x + cos x). No 1º quadrante, f = ($\surd$2-1)(sen x+cos x) = (2-$\surd$2)$\cdot$sen(x+$\pi$/4) $\in$ ($\surd$2-1, 2-$\surd$2{]} $\approx$ (0,414; 0,586{]}, intervalo sem inteiros. No 3º quadrante, f = -($\surd$2+1)(sen x+cos x) = -(2+$\surd$2)$\cdot$sen(x+$\pi$/4) $\in$ ($\surd$2+1, 2+$\surd$2{]} $\approx$ (2,414; 3,414{]}, cujo único inteiro é 3. Conferi com uma amostragem numérica (200 mil pontos em cada quadrante): o 1º quadrante deu de 0,41421 a 0,58579 e o 3º de 2,41421 a 3,41421. Também testei o x dado na resolução, x = 3$\pi$/4 + arcsen(3(2-$\surd$2)/2) $\approx$ 3,429 rad. Ele está em ($\pi$, 3$\pi$/2) e dá f(x) = 3,000. A resolução e o gabarito chegam a k = 3 pelo mesmo caminho.%
\end{itemize}%
\paragraph{Pontos de dúvida}%
\begin{itemize}%
\item%
\textit{(key)} Na simplificação inicial, o gabarito escreve f(x) = $\surd$2 cos x + $\surd$2 sen x - sen x - cos x, sem os módulos (o certo é $\surd$(2cos$^2$x) = $\surd$2|cos x|). Depois ele corrige isso ao tratar cada quadrante com o sinal certo, então o resultado não muda. Mesmo assim, como passagem isolada, a igualdade está errada.%
\item%
\textit{(key)} O gabarito mistura desigualdades estritas e não estritas no 1º quadrante: escreve $\surd$2/2 $\leq$ sen(x+$\pi$/4) $\leq$ 1, depois $\surd$2/2 < f/(2-$\surd$2) $\leq$ 1 e por fim $\surd$2-1 $\leq$ f $\leq$ 2-$\surd$2. Além disso, usa a mesma letra k para o período (2k$\pi$) e para o inteiro procurado. São só questões de redação e não afetam a resposta.%
\item%
\textit{(statement)} O enunciado não diz se os pontos sobre os eixos (x = 0, $\pi$/2, $\pi$, 3$\pi$/2) contam como parte dos quadrantes. A resolução trata esse caso no passo 5: nesses pontos f vale $\surd$2-1 ou $\surd$2+1, que não são inteiros. Por isso a resposta k = 3 vale com qualquer das duas convenções.%
\end{itemize}%
\end{document}